\documentclass{situationist-preprint}

\title{Brouwer's Arguments Reconsidered}
\msid{MTH-07-b}
\msfamily{MTH-07}
\msdiscipline{Philosophy of Mathematics}
\mstype{Companion manuscript}
\mskeywords{Brouwer; intuitionism; excluded middle; strong counterexamples; continuity principle; choice sequences}

\begin{document}
\maketitle

\begin{abstract}
MTH-07-a answers Brouwer's argument from illegitimate extrapolation and his weak counterexamples. Two further arguments of Brouwer's are stronger: the strong counterexamples of intuitionistic analysis, which refute universally quantified forms of excluded middle outright, and the constructive conception of truth on which the whole position rests. We show that the strong counterexamples refute no instance of the law, and that they are theorems about a different subject matter, the information available from finite stages of a choice sequence, of which intuitionistic logic is the correct logic. The constructive conception of truth, taken literally, makes mathematical truth tensed and personal, which is a reason to reject it.
\end{abstract}

\section{The strong counterexamples}\label{s:strong}

A weak counterexample shows only that we are not in a position to assert an instance of excluded middle. Brouwer also gave strong ones. In intuitionistic analysis, real numbers are given by choice sequences of rational approximations, which need not be generated by any law, and a function on the reals must compute each value from a finite initial segment of its argument. From this Brouwer derived the continuity theorem: every total function from the unit interval to the reals is continuous \citep{Heyting1956}. The classical function that takes the value $0$ at $0$ and $1$ elsewhere is total if $\forall x\,(x = 0 \vee x \neq 0)$ holds, and is discontinuous; so intuitionistic analysis proves
\[ \neg\, \forall x\, (x = 0 \vee x \neq 0). \]
This is not a lack of evidence. It is a refutation, and it seems to contradict Theorem~3.2 of MTH-07-a, that every instance of excluded middle is true.

\begin{proposition}\label{p:noinst}
The strong counterexamples refute no instance of excluded middle.
\end{proposition}

\begin{proof}
The refuted sentence is a universal generalization of instances, not an instance. By Proposition~3.1 of MTH-07-a, proved intuitionistically, the negation of every instance is refutable, so that $\forall x\, \neg\neg(x = 0 \vee x \neq 0)$ is an intuitionistic theorem. Intuitionistic logic does not license the step from this to $\forall x\,(x = 0 \vee x \neq 0)$, and the strong counterexample exploits exactly that gap.
\end{proof}

The question is then what the refuted generalization means for the intuitionist. Its proof uses the continuity principle, and the continuity principle is classically false. It is true, however, on a reading the classical mathematician can state: a property of choice sequences is decided by a procedure only if the procedure decides it on the basis of a finite initial segment. Read so, the intuitionistic sentence $\forall x\,(x = 0 \vee x \neq 0)$ says that a procedure can decide, from finite information about any $x$, whether $x$ is zero, and that is false: no finite segment of approximations to $0$ excludes that later approximations move away from it. The strong counterexample is a correct theorem about decidability from finite information. LOG-02 holds that classical logic is the one true logic and that non-classical logics correctly describe other relations; the logic of what is decidable from finite information is one of them, and the strong counterexamples belong to it. They do not bear on whether every real number is or is not zero.

\section{The constructive conception of truth}\label{s:truth}

Behind both kinds of counterexample stands Brouwer's view that mathematics is a languageless construction of the mind, and that a mathematical statement is true only when a construction verifying it has been carried out \citep{Brouwer1907,Brouwer1908}. If that were so, the law would indeed fail, since for an undecided $A$ no construction has verified $A$ and none has verified $\neg A$. Brouwer's later creating-subject arguments make the view explicit by reasoning about the stages at which an idealized mathematician has or has not proved a statement \citep{TroelstraVanDalen1988}.

\begin{objection}[Truth is construction]
Classical determinacy presupposes a realm of mathematical objects independent of construction. The intuitionist denies that there is one, and so denies the premise from which Theorem~3.2 of MTH-07-a follows.
\end{objection}

\begin{reply}
Taken literally, the constructive conception makes truth tensed and relative to a subject: a statement becomes true when someone carries out a construction, and is true for that subject. The creating-subject arguments embrace this, indexing truth to the stages of an idealized mathematician. But a sentence of arithmetic does not say anything about stages or subjects, and if it is proved tomorrow, what is proved is that it was true today. Mathematicians who prove a theorem independently prove the same thing, and the convergence is one of the marks of discovery that MTH-03 uses. The independent realm the objection mentions is, for arithmetic, the structure fixed up to isomorphism by Dedekind's categoricity theorem, on which MTH-07-a rests determinacy; the intuitionist must deny either that the theorem holds or that it says what it does, and neither denial is supported by the constructive conception except by assuming it.
\end{reply}

\section{Conclusion}

Brouwer's strong counterexamples are theorems about what can be decided from finite information, and intuitionistic logic is the right logic for that subject. They refute no instance of excluded middle. His conception of truth as construction would make mathematical truth tensed and personal, which it is not. Classical logic remains the correct logic for mathematics.

\begin{thebibliography}{9}\small
\bibitem[Brouwer(1907)]{Brouwer1907} Brouwer, L.~E.~J. (1907). \emph{Over de grondslagen der wiskunde}. Maas and van Suchtelen.
\bibitem[Brouwer(1908)]{Brouwer1908} Brouwer, L.~E.~J. (1908). De onbetrouwbaarheid der logische principes. \emph{Tijdschrift voor Wijsbegeerte}, 2, 152--158.
\bibitem[Heyting(1956)]{Heyting1956} Heyting, A. (1956). \emph{Intuitionism: An Introduction}. North-Holland.
\bibitem[Troelstra and van Dalen(1988)]{TroelstraVanDalen1988} Troelstra, A.~S. and van Dalen, D. (1988). \emph{Constructivism in Mathematics: An Introduction}, vol.~1. North-Holland.
\end{thebibliography}

\end{document}
